Un fil de cuivre de \qty{\v{d}}{mm} de diamètre transporte un courant de \qty{\v{I}}{A}. Le cuivre compte \num{8.49e28} électrons libres par mètre cube ($e = \qty{1.602e-19}{C}$).

\begin{enumerate}[label=\alph*)]
    \item Déterminez la densité de courant dans le fil.
    \item Déterminez la vitesse de dérive des électrons libres.
    \item Combien de temps faut-il en moyenne à un électron pour parcourir \qty{1}{m} de fil ? Pourquoi la lampe s'allume-t-elle pourtant presque instantanément quand on ferme l'interrupteur ?
\end{enumerate}
